\subsection{平面の三次式からの垂直型の量子化}
\label{subsec:quant-plane}

三次式 $f\in\mathbb C[X_0,X_1,X_2]_3$ に対応する
$\mathbb P^2$ の Poisson 構造を $\sigma_f$ と書く。
アフィン座標 $u=X_1/X_0$、$v=X_2/X_0$ では
$\sigma_f=f(1,u,v)\partial_u\wedge\partial_v$ と規格化する。
二次式の対称化を $g\mapsto\widehat g$ とし、
\[
 A_{f,t}=\mathbb C[t]\langle X_0,X_1,X_2\rangle\Big/
 \left(
 [X_0,X_1]-\frac t3\widehat{f_{X_2}},\quad
 [X_1,X_2]-\frac t3\widehat{f_{X_0}},\quad
 [X_2,X_0]-\frac t3\widehat{f_{X_1}}
 \right)
 \tag{E1}
\]
と置く。Mori--Smith \cite[\S1.7.3、定理1.4、\S1.8.2]{MoriSmith}
の構成と正則性判定により、$0\in U\subset\mathbb A^1_t$ なる
Zariski 開集合上で、これは次数ごとに平坦な三次元 quadratic AS-regular
代数の族である。各次数 $n$ の階数は $\binom{n+2}{2}$ である。
$f=0$ および特異な零三次曲線も許される。

正則性の開性はここでは有限の条件で判定できる。
三つの関係式は、交代三次テンソルに $-(t/3)\widehat f$ を加えた
superpotential の巡回偏微分である。
関係式の独立性を保つ開集合では標準性を満たし、AS 正則性は
対応する $3\times3$ 行列の $2\times2$ 小行列式が
$\mathbb P^2$ に共通零点を持たないことと同値である。
この条件は射影性から開条件であり、$t=0$ の多項式環で成立する。
各次数の商の次元がこの開集合上で一定なので、各斉次部分は局所自由である。

関係式から、アフィン錐上の半古典極限は
\[
 \{X_0,X_1\}=\frac13 f_{X_2},\qquad
 \{X_1,X_2\}=\frac13 f_{X_0},\qquad
 \{X_2,X_0\}=\frac13 f_{X_1}.
\]
Euler の恒等式を使うと
\[
 \{u,v\}
 =\frac{X_0 f_{X_0}+X_1f_{X_1}+X_2f_{X_2}}{3X_0^3}
 =f(1,u,v).
\]
従って $\operatorname{qgr}A_{f,t}$ は指定した $\sigma_f$ を量子化する。
さらに $P=\mathcal O_U[Y_0,Y_1]$ との Segre 積
\[
 C_{f,t}=A_{f,t}\mathbin{\circ}P
   =\bigoplus_{n\geq0}(A_{f,t})_n\otimes_{\mathcal O_U}P_n
 \tag{E2}
\]
を取ると、$t=0$ は $\mathbb P^2\times\mathbb P^1$ の
$\mathcal O(1,1)$ に関する切断環である。
半古典極限の混合括弧は零であり、垂直型成分の任意の元を得る。
次数 $n$ の階数は $(n+1)\binom{n+2}{2}$ である。

この積の圏は、$\mathbb P^1_U$ 上の定数族 $\operatorname{qgr}A_{f,t}$
としても記述できる。アフィン開集合ごとには中央多項式拡大を取り、
その間を通常の $\mathbb P^1$ の座標変換で貼り合わせる。
各因子の ample sequence の外積から、同じ圏が上の Segre 切断環で復元される。
Segre 積と二重次数付き商圏の同値については
\cite[\S3.2、定理3.32]{BallardFarman} を用いる。

\subsection{双二次式からの二因子型の量子化}
\label{subsec:quant-quadric}

二次元ベクトル空間 $A,B$ を固定する。
$\mathbb P(A^*)\times\mathbb P(B^*)$ の Poisson 構造は、
体積形式を選ぶと $D\in\operatorname{Sym}^2 A\otimes\operatorname{Sym}^2 B$
で表される。$A$ の基底を $a_0,a_1$、$B$ の基底を $b_0,b_1$ とし、
$A\otimes B\otimes A\otimes B$ において
\begin{align*}
 w_0={}&a_0b_0a_1b_1-a_1b_0a_0b_1\\
       &-a_0b_1a_1b_0+a_1b_1a_0b_0
\end{align*}
と置く。ここで並置はテンソル積を意味する。
$D$ を第一・第三因子および第二・第四因子に関して対称化した
テンソルを $\widetilde D$ と書き、
\[
                   w(t)=w_0-t\widetilde D
                   \tag{E3}
\]
と定める。

$V_i=A$ または $B$ を $i$ の偶奇で定め、$w_i(t)$ を $w(t)$ の
$i$ 回の巡回移動とする。$V_{i+3}^*$ による $w_i(t)$ の縮約の像を
三次の関係式とし、これらでテンソル $\mathbb Z$-代数を割ったものを
$\mathcal A_{D,t}$ と書く。
Van den Bergh の geometric quintuple の判定と分類定理
\cite[\S5.2、定理5.2.2、定義5.3.1、定理5.5.10]{VdBQuadrics} により、
$0$ を含む Zariski 開集合上で全てのファイバーが
三次元 cubic AS-regular $\mathbb Z$-代数になる。
実際 geometric 条件の失敗は、隣接二因子の非零共変ベクトルによる
縮約が零になるという射影的 incidence の像であり閉条件である。
$w_0$ は linear quadric のテンソルなのでこの閉集合の外にある。
定理によって全ファイバーの各斉次部分の次元は一定であり、
関係式の行列の階数も一定になるため、族は次数ごとに平坦である。

$w(t)$ は二回の巡回移動で不変なので、$\mathcal A_{D,t}$ は
選んだ $V_{i+2}=V_i$ によって二周期的である。
従って偶数 Veronese
\[
 B_{D,t}=\bigoplus_{n\geq0}(\mathcal A_{D,t})_{0,2n}
 \tag{E4}
\]
には通常の $\mathbb N$-次数付き代数構造が入り、
\[
 \dim(B_{D,t})_n=(n+1)^2,\qquad
 B_{D,0}=\bigoplus_{n\geq0}
 H^0(\mathbb P^1\times\mathbb P^1,\mathcal O(n,n)).
\]
この Veronese は同じ非可換 quadric の圏を与える。

この族の Kodaira--Spencer 写像は $D$ に対応する Poisson 構造である。
ここではその規格化も確かめる。
写像は $\operatorname{SL}(A)\times\operatorname{SL}(B)$-同変であり、
始域および
\[
 HH^2(\mathbb P^1\times\mathbb P^1)
 =H^0(\mathcal O(2,2))
 =\operatorname{Sym}^2A\otimes\operatorname{Sym}^2B
\]
は同じ既約表現である。従って最高ウェイトベクトル
$D=a_0^2b_0^2$ で確認すればよい。
全ての位置の $a_0,b_0$ に対応する矢を universal localization で可逆にし、
それらによって各対象を同定する。
各位置に残る矢を $u_k,v_k$ とすると、三次関係式のうち二つは
$u_{k+1}=u_k$ および $v_{k+1}=v_k$ を与える。
残る関係式は、共通の座標 $u,v$ に対し
\[
                         uv-vu=t
\]
となる。この局所表示は Weyl 代数の族であり、PBW 基底 $u^iv^j$ を持つ。
古典的な点では $u=a_1/a_0$、$v=b_1/b_0$ なので、半古典極限は
$D(1,u,1,v)\partial_u\wedge\partial_v$ である。

以上から
\[
 C_{D,t}=B_{D,t}\mathbin{\circ}\mathcal O_U[Y_0,Y_1]
 \tag{E5}
\]
は三重積の任意の二因子型 Poisson 構造を量子化する代数的な族である。
各次数の階数は $(n+1)^3$、$t=0$ の切断環は
$\mathcal O(1,1,1)$ に関する Segre 環である。
$D$ の零曲線の滑らかさは仮定しない。

\subsection{同変性と対数型への適用}
\label{subsec:quant-elliptic-equivariance}

式 (E3) は二因子の線形座標変換と両立する。
特に代数群が $w_0$ と $\widetilde D$ を同じ指標で変換するなら、
全ての $B_{D,t}$ にその群の次数付き作用がある。
三重積の一般の対数型では、二つの $(1,1)$ 因子の交点を
$(0,0)$ と $(\infty,\infty)$ に送って
\[
 F=f x+g z,\qquad G=h x+k z
\]
と書ける。本書 の対数型の表示を用いると
\[
 \pi=\sigma+y\partial_y\wedge V,\qquad
 \sigma=aFG\partial_z\wedge\partial_x,\qquad
 V=\kappa(x\partial_x+z\partial_z),\quad
 \kappa=b(fk-gh)
 \tag{E6}
\]
である。$(u,v)=(x,z)$ に対して $D=-aFG$ と取る。
$\sigma$ は radial $\mathbb G_m$ 作用で不変である。
斉次座標の重みを $(0,1)$ と $(0,1)$ に取ると、
$w_0$ と $\widetilde D$ はともに重み $2$ を持つため、
上の量子化は同じ $\mathbb G_m$ 作用 $\rho$ を保つ。

$q=1+t\kappa$ を可逆にした開集合に制限し、
\[
 T=B_{D,t}[Y_0][Y_1;\rho_q],\qquad
 Y_1c=\rho_q(c)Y_1,\qquad Y_0\text{ は中央}
 \tag{E7}
\]
という二重次数付き Ore 拡大を作る。
次数は $\deg c=(\deg_B c,0)$、$\deg Y_i=(0,1)$ とする。
その対角部分 $\bigoplus_{n\geq0}T_{(n,n)}$ は次数ごとに平坦であり、
各次数の階数は $(n+1)^3$ である。
アフィン座標 $y=Y_1/Y_0$ について
\[
 \{y,c\}=\kappa y(x\partial_x+z\partial_z)(c)
\]
なので、半古典極限は式 (E6) と一致する。
従って各対数型成分の一般元についても、代数的な
Segre 切断環の族が得られる。

\subsection{pencil 成分のパラメーターと非可換 quadric}
\label{subsec:quant-pencil-parameters}

三重積の $(1,1,1)$-pencil の基底楕円曲線 $E$ は、
三つの射影から次数 $2$ の直線束 $L_0,L_1,L_2$ を持つ。
一般の場合、この四つ組は Van den Bergh の elliptic quadruple の条件を満たす。
また pencil の二つの生成元を第四の二次元因子にまとめると、
geometric quintuple の四重テンソルが得られる。
これは、本書 の pencil 成分の射影 GIT 商
$\mathbb P(1,2,2,3)$ と、Okawa--Ueda
\cite[定理1.1]{OkawaUedaQuadrics} の非可換 quadric のコンパクトモジュライ
$\mathbb P(2,4,4,6)$ が同じ重み付き射影空間になることを説明する。

この四重テンソルは非可換 quadric 曲面を与える。
本書 の「最大 pencil 成分の Segre 射影量子化」では、
$L_E=L_0\otimes L_1\otimes L_2$ に対する filtered deformation の
相対 $\mathbb P^1$ ruling から三次元の斉次環を構成する。
二つの次数一の中央元 $s,t$ を持ち、各次数の階数は $(n+1)^3$ である。
さらに同論文の「楕円曲線を中心とする基本変換の相対共通圏」では、
楕円曲線のグラフに沿う gluing と二回の mutation により、
元の三重積の量子化を一方の収縮とする二つの SOD を得る。
